% `tex.web` §846 narrates the line breaker, and the narration is a faithful
% picture of the active list — so getting it right is getting §829 right.
%
% §835 walks the list forward and creates a line-number group's nodes the
% moment the cursor crosses into the next group:
%
%     begin l:=line_number(r);
%     if l>old_l then
%       begin if (minimum_demerits<awful_bad)and
%          ((old_l<>easy_line)or(r=last_active)) then
%         @<Create new active nodes for the best feasible breaks just found@>;
%
% so the log runs candidates, nodes, candidates, nodes — not every candidate
% of a breakpoint and then every node of it.
%
% §845 links each new node in before the node the cursor is standing on, and
% runs `prev_r:=q` after each fitness class:
%
%     link(q):=r; link(prev_r):=q; prev_r:=q;
%
% so a group's four classes go in one after another in class order, and the
% group as a whole goes *in front of* the run it joins — the cursor is already
% on the first node of the next group when the flush happens. The merged group
% is the exception: it is flushed at `r=last_active`, so it goes on the end.
%
% §857 prints the material since the last breakpoint reported, and prints
% nothing when there is none — `print_nl("")` starts no line when none is in
% progress. We pushed an empty line, which put a blank between the last `@@`
% node of a paragraph and its `\par` breaks.
%
% `\parshape` is what makes the groups distinguishable: without it every line
% has the same measure, `easy_line` is nought, and the whole paragraph is one
% merged group.
\catcode`\{=1 \catcode`\}=2
\font\r=cmr10 \r \hsize120pt \parindent20pt \baselineskip12pt \parskip6pt
\tolerance=10000 \pretolerance=100 \linepenalty=10 \adjdemerits=10000
\hbadness=10000 \vbadness=10000
\parfillskip=0pt plus 1fil
\rightskip 0pt plus 1fil
\parshape=4 0pt 120pt 10pt 100pt 20pt 80pt 0pt 120pt
\tracingparagraphs1 \tracingonline0
\noindent the quick brown fox jumps over the lazy dog, while office traffic waffles fluffily past the coffee shop on the corner of the avenue near the old railway station.
\par
\tracingparagraphs0
\end
